%%%%%%%%%%%%%%%%%%%%%%%%%%%%%%%%%%%%%%%%%%%%%%%%%%%%%%%%%%%%%%%%%%%%%%%
% Topological Modulation of Transient Dimensional Inflation
% in Energy-Coupled Chemical Oscillators
%
% Paper 5 — March 2026
%%%%%%%%%%%%%%%%%%%%%%%%%%%%%%%%%%%%%%%%%%%%%%%%%%%%%%%%%%%%%%%%%%%%%%%
\documentclass[twocolumn,showpacs,preprintnumbers,amsmath,amssymb,prd,floatfix]{revtex4-2}
\usepackage{graphicx}
\usepackage{amsmath}
\usepackage{amssymb}
\usepackage{bm}
\usepackage{hyperref}
\usepackage{xcolor}
\usepackage{booktabs}
\usepackage{multirow}

\newcommand{\VV}[1]{{\color{red}{#1}}}
\newcommand{\tauexp}{\tau_{>1.2}}
\newcommand{\Dtwo}{D_2}
\newcommand{\kcat}{k_{\mathrm{cat}}}
\newcommand{\Kd}{K_d}
\newcommand{\Gtot}{G_{\mathrm{total}}}
\newcommand{\Rw}{R_w}
\renewcommand{\sd}{\sigma_d}

\begin{document}

\title{Topological Modulation of Transient Dimensional Inflation\\
in Energy-Coupled Chemical Oscillators}

\author{Christopher Merrill}
\affiliation{Independent Researcher}
\author{With computational collaboration from Claude (Anthropic), ChatGPT (OpenAI), and Gemini (Google)}

\date{\today}

\begin{abstract}
Papers~1--4 in this series established that timescale separation between fast oscillatory cores and a slow shared energy reservoir transiently inflates the correlation dimension $\Dtwo$ beyond the limit-cycle floor, and that evolutionary selection can prolong this exploratory phase. Paper~4 showed that expanding the mutable parameter space from one to three dimensions accelerated neutral decay of transient retention, raising the question: can reaction-network topology buffer this fragility? Here we test this directly. We augment the enzyme-complex coupled Brusselator with five designed motifs and 20~random-wiring controls, all coupled through the shared energy pool via a standardized embedding protocol. A ridge-width screen ($N = 1{,}800$~parameter samples per topology, 14{,}400 total) reveals that all designed motifs \textit{suppress} transient dimensionality relative to the unaugmented baseline ($\Rw = 0.04$--$0.09$ vs.\ $0.37$; Mann--Whitney $p < 10^{-97}$). The 20~random wirings split bimodally: 8~are completely dead ($\Rw = 0$), while the surviving 12~have $\Rw = 0.02$--$0.40$, with two exceeding the baseline. A drift-sensitivity test across seven perturbation magnitudes ($\sd = 0$--$0.5$, 1{,}202~evaluations per topology) finds that all topologies degrade at the same rate ($\sim$42\%~$\Rw$ loss at $\sd = 0.5$), indicating that topology controls the size of the viable parameter region but not its sensitivity to drift. The alive/dead split among random wirings is predicted by a single graph-theoretic feature: whether the motif species catalyzing $E \to E_w$ is a net concentration accumulator in the motif subgraph (Fisher's exact $p = 0.0007$; 0/12~alive vs.\ 6/8~dead wirings are accumulators). This structural predictor also explains why all designed motifs fail: the embedding protocol routes energy coupling through a cycle endpoint that necessarily accumulates concentration. We conclude that network augmentation generically destroys transient complexity via energy drain, and that the rare wirings preserving it satisfy a simple structural constraint: the energy-coupling point must disperse rather than accumulate mass flow.
\end{abstract}

\maketitle

%%%%%%%%%%%%%%%%%%%%%%%%%%%%%%%%%%%%%%%%%%%%%%%%%%%%%%%%%%%%%%%%%%%%%%%
\section{Introduction}
\label{sec:introduction}
%%%%%%%%%%%%%%%%%%%%%%%%%%%%%%%%%%%%%%%%%%%%%%%%%%%%%%%%%%%%%%%%%%%%%%%

The first four papers in this series traced a progression from dynamical rigidity through transient escape to evolutionary tuning of that escape. Paper~1~\cite{paper1} showed that random autocatalytic growth dilutes the activation ratio $\eta = \Dtwo/r_S$. Paper~2~\cite{paper2} demonstrated that feedback-aligned growth preserves viability but cannot lift $\Dtwo$ above the limit-cycle floor. Paper~3~\cite{paper3} identified the mechanism that breaks this floor: timescale separation between fast oscillatory cores and a slow shared energy reservoir transiently inflates $\Dtwo$ to $\sim$1.7 for hundreds of oscillator periods. Paper~4~\cite{paper4} showed that tournament selection on the transient retention metric $\tauexp$ (the count of 2{,}500-unit windows with $\Dtwo > 1.2$) drives a $2.8\times$ increase in mean $\tauexp$ via directional convergence of the energy dissipation rate $\gamma$.

Paper~4 also identified a vulnerability. When the mutable parameter space was expanded from one dimension ($\gamma$ alone) to three ($\gamma, J, \kcat$), neutral drift collapsed $\tauexp$ from $\sim$3.1 to 0.55 over 40~generations, with 12.9~lethal phenotypes per generation. The viable region of parameter space---the ``evolutionary ridge''---appears narrow relative to the mutation step size. This raises a natural question: \textit{can reaction-network topology widen this ridge?}

The hypothesis has intuitive appeal. Network motifs that create redundant slow-variable pathways might provide alternative routes to timescale separation, expanding the parameter-space volume that sustains transient chaos. In the language of chemical reaction network theory (CRNT), additional complexes and linkage classes can increase the deficiency of a network~\cite{feinberg1972, horn1972}, potentially enriching its dynamical repertoire.

We test this hypothesis with a systematic topology screen. Five designed motifs (simple cycle, parallel dual-path, branch-rejoin, feedforward-return, nested dual loop) and 20~random-wiring controls are grafted onto the Paper~3 base model via a standardized embedding protocol that couples all augmentations through the shared energy pool~$E$. We measure the ``ridge width'' $\Rw$---the fraction of random parameter draws sustaining $\tauexp > 0$---as the primary metric of topological buffering.

The results reject the original hypothesis but reveal a more nuanced story. All designed motifs suppress transient dimensionality. Random wirings split bimodally into alive and dead categories. And the alive/dead distinction is predicted by a single structural property of the motif's relationship to the energy-coupling point.

%%%%%%%%%%%%%%%%%%%%%%%%%%%%%%%%%%%%%%%%%%%%%%%%%%%%%%%%%%%%%%%%%%%%%%%
\section{Methods}
\label{sec:methods}
%%%%%%%%%%%%%%%%%%%%%%%%%%%%%%%%%%%%%%%%%%%%%%%%%%%%%%%%%%%%%%%%%%%%%%%

\subsection{Base Model}

We use the enzyme-complex coupled Brusselator from Papers~3--4: two Brusselator oscillator cores sharing a slow energy pool $E$, gated by an enzyme complex $G/GE$. The model comprises 14~reactions with mass-action kinetics. Dynamic species tracked for $\Dtwo$: $\{X_1, Y_1, X_2, Y_2, E\}$ (5~variables). Integration via LSODA over $t \in [0, 20{,}000]$ with 40{,}000 time points, rtol $= 10^{-6}$, atol $= 10^{-9}$, max\_step $= 10.0$. Fixed parameters: $A = 1.0$, $B = 3.0$, $k_\text{on} = k_\text{off} = 10.0$, $\Gtot = 1.0$. See Paper~3 for the full reaction scheme.

\subsection{Topology Panel}

We construct eight topologies organized into two groups (Table~\ref{tab:topology_panel}):

\begin{itemize}
\item \textbf{Group~A (controls):} T0 (unaugmented baseline), T1a and T1b (random-wiring controls with specific seeds, matched for species and reaction count to the largest Group~B motif).

\item \textbf{Group~B (designed motifs):} T2 (simple 3-cycle, $A_m \to B_m \to C_m \to A_m$), T3 (parallel dual-path cycle), T4 (branch-and-rejoin with asymmetric weights), T5 (feedforward with return), T6 (nested dual loop with inner 2-cycle).
\end{itemize}

\begin{table}[b]
\centering
\caption{Topology panel. All Group~B motifs and random controls add species and reactions to the 14-reaction, 12-species base model. Cyclo: cyclomatic number. Def: deficiency.}
\label{tab:topology_panel}
\begin{tabular}{llcccc}
\toprule
Name & Description & Species & Rxns & Cyclo & Def \\
\midrule
T0 & Baseline & 12 & 14 & 17 & 2 \\
T1a & Random ctrl (s=1001) & 16 & 20 & 20 & 3 \\
T1b & Random ctrl (s=2002) & 16 & 20 & 19 & 3 \\
T2 & Simple 3-cycle & 15 & 18 & 19 & 3 \\
T3 & Parallel dual-path & 16 & 20 & 20 & 3 \\
T4 & Branch-rejoin & 16 & 20 & 20 & 3 \\
T5 & Feedforward-return & 15 & 19 & 20 & 3 \\
T6 & Nested dual loop & 16 & 20 & 20 & 3 \\
\bottomrule
\end{tabular}
\end{table}

\subsection{Embedding Protocol}
\label{sec:embedding}

All topological augmentations follow a standardized embedding protocol to ensure fair comparison:

\begin{enumerate}
\item New species ($R_{m0}, R_{m1}, \ldots$) are initialized at concentration 1.0.
\item Intra-motif reactions follow mass-action kinetics with rate constants sampled from log-uniform$(0.1, 10)$.
\item Each motif couples to the base system through exactly one energy-drain reaction of the form $R_{mk} + E \to R_{mk} + E_w$, where $R_{mk}$ is a designated ``coupler'' species within the motif. This catalytic drain is the sole interface between the motif and the base oscillator.
\item No motif species participates directly in any base-model reaction (no coupling to $X_i$, $Y_i$, $G$, or $GE$).
\item The $\Dtwo$ projection includes all motif species except waste ($W_i$), chemostat ($A$, $B$, $E_\text{src}$), and gate ($G$, $GE$) species.
\end{enumerate}

\subsection{Random-Wiring Ensemble}
\label{sec:random_ensemble}

Twenty random-wiring controls are generated to test whether any observed topological effects exceed what is expected from randomly structured augmentations with matched complexity. Each random wiring adds 4~species and 6~reactions (matching T3/T4/T6), with 5~intra-motif conversion reactions drawn uniformly among the new species (no self-loops) and 1~energy-coupling reaction to a randomly selected motif species. The 20~wirings use deterministic pseudo-random seeds for reproducibility.

\subsection{Ridge-Width Screen (Phase~1)}
\label{sec:phase1}

The primary metric is the ridge width $\Rw$: the fraction of random parameter draws that yield $\tauexp > 0$. For each named topology, we sample:

\begin{itemize}
\item 300 base-parameter draws: $\gamma \in [0.0005, 0.01]$ log-uniform, $J \in [2, 12]$ uniform, $\kcat \in [0.05, 0.8]$ uniform.
\item 3 independent motif-rate draws per base draw (log-uniform in $[0.1, 10]$ for each motif reaction).
\item 2 integration seeds per configuration (seeds 42, 179).
\end{itemize}

Total: $300 \times 3 \times 2 = 1{,}800$ evaluations per named topology, 14{,}400 across the 8-topology panel. For the extended random ensemble, each of the 20~rewirings receives 100~base-parameter draws $\times$ 2~seeds $= 200$ evaluations, totaling 4{,}000 additional evaluations. The screen ran on 10~AWS EC2 spot instances (t3.micro, us-west-2) over approximately 45~hours.

Each evaluation has a 120\,s wall-clock timeout; timed-out evaluations receive $\tauexp = 0$.

\subsection{Drift-Sensitivity Test (Phase~2)}
\label{sec:phase2}

To test whether topology buffers sensitivity to parameter perturbation, we measure how $\Rw$ degrades under Gaussian drift. Four candidate topologies are selected based on Phase~1 results (T0, T6, rw16, rw4). For each, we identify a ``center point''---the parameter combination with the highest mean $\tauexp$ from Phase~1---and evaluate $\Rw$ at seven drift magnitudes:

\begin{equation}
\sd \in \{0, 0.05, 0.1, 0.15, 0.2, 0.3, 0.5\}
\end{equation}

At each $\sd > 0$, 100~draws are sampled from $\mathcal{N}(\theta_\text{center}, \sd^2 \cdot \text{diag}(\theta_\text{center}^2))$ in each parameter dimension, with 2~seeds per draw (200~evaluations per drift level). At $\sd = 0$, a single draw (2~evaluations) suffices since the result is deterministic. Total: 1{,}202 evaluations per topology, 4{,}808 across 4~candidates. This phase ran on 4~spot instances over approximately 30~hours.

\subsection{Structural Predictor Analysis (Phase~4)}
\label{sec:phase4}

For each of the 20~random wirings, we reconstruct the motif subgraph (the directed graph among $\{R_{m0}, R_{m1}, R_{m2}, R_{m3}\}$, excluding the $E$-coupling reaction) and compute:

\begin{itemize}
\item \textbf{Net flow into the E-coupler:} $f_\text{net} = d_\text{in} - d_\text{out}$, where $d_\text{in}$ and $d_\text{out}$ are the in-degree and out-degree of the E-coupled species within the motif subgraph.
\item \textbf{Terminal sink:} whether the E-coupler has $d_\text{out} = 0$ (no outgoing motif edges).
\end{itemize}

The association between structural features and the alive/dead classification ($\Rw > 0$ vs.\ $\Rw = 0$) is tested using Fisher's exact test.

%%%%%%%%%%%%%%%%%%%%%%%%%%%%%%%%%%%%%%%%%%%%%%%%%%%%%%%%%%%%%%%%%%%%%%%
\section{Results}
\label{sec:results}
%%%%%%%%%%%%%%%%%%%%%%%%%%%%%%%%%%%%%%%%%%%%%%%%%%%%%%%%%%%%%%%%%%%%%%%

\subsection{Phase~1: Designed Motifs Suppress Transient Dimensionality}

All five designed motifs reduce $\Rw$ relative to the unaugmented baseline (Table~\ref{tab:phase1_results}). The effect is large: T0 has $\Rw = 0.372$, while the best designed motif (T6, nested dual loop) achieves only $\Rw = 0.091$. The weakest motifs (T3, T4) are nearly dead at $\Rw = 0.041$. Mann--Whitney $U$ tests confirm that T0~$>$~every Group~B motif ($p < 10^{-97}$ for all comparisons).

\begin{table}[t]
\centering
\caption{Phase~1 ridge-width screen results. $\Rw$: fraction of 1{,}800 evaluations with $\tauexp > 0$. Mean $\tauexp$: average across all evaluations. All Group~B motifs are significantly below T0.}
\label{tab:phase1_results}
\begin{tabular}{llcc}
\toprule
Topology & Group & $\Rw$ & Mean $\tauexp$ \\
\midrule
T0 & A (baseline) & 0.372 & 0.671 \\
T1a & A (random ctrl) & 0.003 & 0.003 \\
T1b & A (random ctrl) & 0.335 & 0.588 \\
\midrule
T2 & B (3-cycle) & 0.062 & 0.065 \\
T3 & B (dual-path) & 0.041 & 0.045 \\
T4 & B (branch-rejoin) & 0.041 & 0.045 \\
T5 & B (feedforward) & 0.053 & 0.059 \\
T6 & B (nested loop) & 0.091 & 0.102 \\
\bottomrule
\end{tabular}
\end{table}

The two Group~A random controls are themselves informative: T1a is nearly dead ($\Rw = 0.003$) while T1b is close to baseline ($\Rw = 0.335$). This foreshadows the bimodality observed in the full random ensemble.

\begin{figure}[t]
\centering
\includegraphics[width=0.95\columnwidth]{figures/fig1_rw_distribution.png}
\caption{Phase~1 ridge-width screen. Blue bars: Group~A controls (T0 baseline, T1a/T1b random). Red bars: Group~B designed motifs. Dashed green line and shaded band: random-wiring ensemble mean $\pm 1\sigma$. All designed motifs fall below both the baseline and the random ensemble mean.}
\label{fig:rw_distribution}
\end{figure}

\subsection{Random-Wiring Ensemble: Bimodal Viability}

The 20~random wirings divide sharply into two classes (Fig.~\ref{fig:rw_distribution}). Eight wirings are completely dead ($\Rw = 0$), while the remaining twelve have $\Rw$ ranging from 0.020 to 0.400 (Table~\ref{tab:random_ensemble}). Two random wirings (rw16 and rw17) exceed the T0 baseline, with rw16 reaching $\Rw = 0.400$.

\begin{table}[b]
\centering
\caption{Random-wiring ensemble results, sorted by $\Rw$. Eight wirings are completely dead; twelve sustain transient dimensionality, with two exceeding the unaugmented baseline T0 ($\Rw = 0.372$).}
\label{tab:random_ensemble}
\begin{tabular}{lcccl}
\toprule
Wiring & $\Rw$ & Mean $\tauexp$ & $f_\text{net}$ & Class \\
\midrule
rw16 & 0.400 & 0.740 & $-2$ & Alive \\
rw17 & 0.365 & 0.665 & $0$ & Alive \\
rw4 & 0.350 & 0.715 & $-1$ & Alive \\
rw3 & 0.325 & 0.600 & $-2$ & Alive \\
rw1 & 0.320 & 0.560 & $0$ & Alive \\
rw19 & 0.320 & 0.530 & $-4$ & Alive \\
rw2 & 0.310 & 0.440 & $0$ & Alive \\
rw8 & 0.310 & 0.555 & $-1$ & Alive \\
rw13 & 0.275 & 0.505 & $0$ & Alive \\
rw5 & 0.270 & 0.450 & $-1$ & Alive \\
rw7 & 0.260 & 0.450 & $-1$ & Alive \\
rw18 & 0.020 & 0.020 & $-1$ & Alive \\
\midrule
rw0 & 0.000 & 0.000 & $+1$ & Dead \\
rw6 & 0.000 & 0.000 & $+2$ & Dead \\
rw9 & 0.000 & 0.000 & $-1$ & Dead \\
rw10 & 0.000 & 0.000 & $+2$ & Dead \\
rw11 & 0.000 & 0.000 & $+1$ & Dead \\
rw12 & 0.000 & 0.000 & $+1$ & Dead \\
rw14 & 0.000 & 0.000 & $+2$ & Dead \\
rw15 & 0.000 & 0.000 & $-3$ & Dead \\
\bottomrule
\end{tabular}
\end{table}

The mean $\Rw$ across surviving wirings is $0.294 \pm 0.091$, comparable to the T0 baseline. The overall ensemble mean is $0.176 \pm 0.160$, reflecting the bimodal structure (Fig.~\ref{fig:random_bimodality}).

\begin{figure}[t]
\centering
\includegraphics[width=0.95\columnwidth]{figures/fig4_random_bimodality.png}
\caption{Distribution of $\Rw$ across 20~random wirings. The bimodal split is stark: 8~wirings are completely dead ($\Rw = 0$) while 12~cluster around $\Rw \approx 0.3$, with nothing in between. Dashed line: T0 baseline ($\Rw = 0.372$).}
\label{fig:random_bimodality}
\end{figure}

\subsection{Phase~2: Drift Sensitivity Is Topology-Independent}

All four Phase~2 candidates degrade at similar rates under parameter drift (Table~\ref{tab:phase2_results}, Fig.~\ref{fig:drift_curves}). At $\sd = 0.5$, all topologies retain 55--58\% of their undrifted $\Rw$. Linear regression slopes range from $-0.86$ (T0) to $-0.95$ (T6), with no statistically significant difference between topologies.

\begin{table}[t]
\centering
\caption{Phase~2 drift-sensitivity results. $S_{0.5}$: fraction of $\Rw$ retained at $\sd = 0.5$ relative to $\sd = 0.05$. Slope: linear fit of $\Rw$ vs.\ $\sd$.}
\label{tab:phase2_results}
\begin{tabular}{lcccc}
\toprule
Topology & $\Rw(0.05)$ & $\Rw(0.5)$ & $S_{0.5}$ & Slope \\
\midrule
T0 & 1.000 & 0.580 & 0.580 & $-0.86$ \\
T6 & 1.000 & 0.575 & 0.575 & $-0.95$ \\
rw16 & 1.000 & 0.580 & 0.580 & $-0.90$ \\
rw4 & 0.965 & 0.545 & 0.565 & $-0.89$ \\
\bottomrule
\end{tabular}
\end{table}

T6 shows marginally better $\Rw$ retention at moderate drift ($\Rw = 0.870$ at $\sd = 0.2$ vs.\ 0.730 for T0), but this advantage disappears by $\sd = 0.5$. Mean $\tauexp$ decay confirms the same pattern: all topologies retain 42--50\% of their mean $\tauexp$ at $\sd = 0.5$.

\begin{figure}[t]
\centering
\includegraphics[width=0.95\columnwidth]{figures/fig2_drift_curves.png}
\caption{Phase~2 drift sensitivity. $\Rw$ vs.\ perturbation magnitude $\sd$ for four candidate topologies. All four curves track nearly identically, confirming that topology controls ridge \textit{size} but not drift \textit{sensitivity}.}
\label{fig:drift_curves}
\end{figure}

\subsection{Structural Predictor: E-Coupler Accumulation}

We hypothesized that the alive/dead split among random wirings might be predicted by the structural relationship between the motif subgraph and the energy-coupling point. Specifically, the species catalyzing $E \to E_w$ experiences a catalytic E-drain rate proportional to its own concentration. If this species accumulates mass from the motif subgraph (more incoming edges than outgoing), its steady-state concentration rises, increasing the rate of energy drain and starving the $GE$ complex that mediates timescale separation.

We define the net flow $f_\text{net} = d_\text{in} - d_\text{out}$ of the E-coupled species within the motif subgraph. Among the 20~random wirings:

\begin{itemize}
\item \textbf{Alive wirings} ($\Rw > 0$, $n = 12$): 0/12 have $f_\text{net} \geq 1$.
\item \textbf{Dead wirings} ($\Rw = 0$, $n = 8$): 6/8 have $f_\text{net} \geq 1$.
\end{itemize}

Fisher's exact test yields $p = 0.0007$ (Table~\ref{tab:predictor}). The E-coupler's terminal-sink status ($d_\text{out} = 0$) is also predictive but less strongly: 0/12 alive vs.\ 3/8 dead ($p = 0.046$).

\begin{table}[b]
\centering
\caption{Structural predictor of random-wiring viability. $f_\text{net}$: net flow into the E-coupled species from the motif subgraph. Terminal sink: E-coupler has no outgoing motif edges.}
\label{tab:predictor}
\begin{tabular}{lccc}
\toprule
Predictor & Alive & Dead & Fisher $p$ \\
\midrule
$f_\text{net} \geq 1$ (accumulator) & 0/12 & 6/8 & 0.0007 \\
$d_\text{out} = 0$ (terminal sink) & 0/12 & 3/8 & 0.046 \\
\bottomrule
\end{tabular}
\end{table}

Two dead wirings (rw9, rw15) have $f_\text{net} < 1$ but are still dead ($f_\text{net} = -1$ and $-3$, respectively; Fig.~\ref{fig:predictor}). These exceptions suggest that the net-flow predictor captures the dominant mechanism but does not fully determine viability.

\begin{figure}[t]
\centering
\includegraphics[width=0.95\columnwidth]{figures/fig3_predictor.png}
\caption{Structural predictor of random-wiring viability. Each point is one of the 20~random wirings, plotted by E-coupler net flow $f_\text{net}$ (x-axis) and ridge width $\Rw$ (y-axis). Green: alive ($\Rw > 0$). Gray: dead ($\Rw = 0$). Dashed red line: $f_\text{net} = 1$ threshold. No alive wiring has $f_\text{net} \geq 1$; 6/8 dead wirings do (Fisher $p = 0.0007$). Points are jittered horizontally for visibility.}
\label{fig:predictor}
\end{figure} Other factors---such as the specific rate constants sampled, the motif's internal relaxation timescale, or higher-order graph properties---may account for the remaining variance.

\subsection{Why All Designed Motifs Fail}

The structural predictor explains a result that initially appeared surprising: \textit{every} designed motif (T2--T6) suppresses transient dimensionality. In all five designs, the energy-coupling reaction routes through a cycle endpoint ($C_m$ in T2, T5; $D_m$ in T3, T4, T6) that by construction receives net inflow from the motif cycle. This is an artifact of the embedding protocol (\S\ref{sec:embedding}), which routes coupling through a species that acts as a concentration accumulator in all five motif topologies.

The random ensemble, by contrast, assigns the E-coupler uniformly at random among the motif species. When the randomly selected coupler happens to be a net disperser ($f_\text{net} \leq 0$), the wiring survives. This explains both the bimodality and the specific wirings that exceed baseline.

%%%%%%%%%%%%%%%%%%%%%%%%%%%%%%%%%%%%%%%%%%%%%%%%%%%%%%%%%%%%%%%%%%%%%%%
\section{Discussion}
\label{sec:discussion}
%%%%%%%%%%%%%%%%%%%%%%%%%%%%%%%%%%%%%%%%%%%%%%%%%%%%%%%%%%%%%%%%%%%%%%%

\subsection{Topology Controls Ridge Size but Not Drift Sensitivity}

Phases~1 and~2 yield a clean separation. Topology dramatically affects $\Rw$---from 0 (dead wirings) to 0.40 (rw16), a factor of $\infty$---but has no measurable effect on the \textit{rate} of $\Rw$ degradation under drift. All four Phase~2 candidates lose approximately 42\% of their $\Rw$ at $\sd = 0.5$, regardless of their starting $\Rw$ or structural composition.

This suggests that the fragility documented in Paper~4 (rapid neutral decay of $\tauexp$ in the V3b multi-parameter experiment) is an intrinsic property of the underlying dynamical mechanism---the timescale-separation window between the energy pool and the oscillator cores---rather than an artifact of insufficient topological complexity.

\subsection{Network Growth as Dimensional Sink}

The dominant finding of Phase~1 is negative: augmenting the reaction network generically \textit{destroys} transient dimensionality rather than buffering it. The mechanism is straightforward. Every augmentation couples to the base system through a catalytic $E \to E_w$ drain. This additional energy sink lowers the steady-state $[E]$, which reduces $[GE]$ (the enzyme complex mediating slow-fast coupling), which narrows the timescale-separation window that drives transient $\Dtwo$ inflation.

The magnitude of this drain depends on the steady-state concentration of the E-coupled species, which is controlled by the motif's internal flow structure. Species that accumulate concentration (net inflow from other motif nodes) have higher steady-state levels and correspondingly stronger E-drain. This provides the mechanistic link between the graph-theoretic predictor ($f_\text{net}$) and the dynamical outcome ($\Rw$).

\subsection{The Accumulation--Dispersion Axis}

The structural predictor we identify---net flow into the E-coupled species---is simple but specific. It is not a global network property (cyclomatic number, deficiency, and weak reversibility are all identical across the 20~random wirings) but a \textit{local} property of the coupling interface between the motif and the energy pool.

This specificity has a practical implication. If one wished to design a topological augmentation that preserves or enhances transient complexity, the design constraint is clear: the E-coupling point must not be a concentration accumulator. More precisely, the species catalyzing the energy drain must have outgoing edges in the motif subgraph that prevent mass accumulation. This amounts to a flow-balance condition at the coupling interface.

\subsection{Implications for Prebiotic Chemistry}

The results suggest a stringent constraint on prebiotic network growth. If a proto-metabolic system relies on transient dynamical complexity for state-space exploration~\cite{nghe2015, vanchurin2022}, then random addition of reactions is overwhelmingly likely to destroy this capability. Of the 20~random wirings we tested, 8 (40\%) are immediately lethal and most of the survivors show reduced $\Rw$ relative to the unaug\-mented system.

The rare wirings that survive share a structural motif: dispersed energy coupling. In a prebiotic context, this translates to the requirement that new catalytic pathways draining the energy pool must be \textit{leaky}---they must not trap concentration at the drain point. This is reminiscent of the ``parasitic side-reaction'' problem in prebiotic chemistry~\cite{szostak2012}, but viewed from the dynamical rather than the thermodynamic perspective.

Combined with Paper~4's finding that evolutionary selection can tune $\tauexp$ over relatively few generations, the emerging picture is one of a narrow but navigable evolutionary path: random network elaboration almost always destroys transient complexity, but the rare surviving topologies can be further optimized by parameter selection.

\subsection{Limitations}

Several limitations constrain generalization:

\begin{enumerate}
\item The base model is specific (enzyme-complex coupled Brusselator). The extent to which the accumulation--dispersion predictor generalizes to other oscillator architectures is unknown.

\item The sample size (20~random wirings) limits statistical power. The Fisher test is significant but two dead wirings fall outside the predictor ($f_\text{net} = -1$ and $-3$), suggesting additional mechanisms.

\item The embedding protocol constrains all coupling to a single catalytic $E$-drain reaction. Alternative coupling modes (e.g., competitive inhibition, allosteric regulation) might yield different results.

\item We do not vary the number of added species or reactions, only their wiring. The relationship between augmentation \textit{size} and viability is unexplored.
\end{enumerate}

%%%%%%%%%%%%%%%%%%%%%%%%%%%%%%%%%%%%%%%%%%%%%%%%%%%%%%%%%%%%%%%%%%%%%%%
\section{Conclusions}
\label{sec:conclusions}
%%%%%%%%%%%%%%%%%%%%%%%%%%%%%%%%%%%%%%%%%%%%%%%%%%%%%%%%%%%%%%%%%%%%%%%

Three principal conclusions emerge:

\begin{enumerate}
\item \textbf{Network augmentation generically suppresses transient complexity.} All five designed motifs and 40\% of random wirings destroy transient dimensional inflation. The mechanism is energy drain through the coupling interface.

\item \textbf{Topology controls ridge size but not drift sensitivity.} The viable parameter region varies by a factor of $>$10$\times$ across topologies, but all topologies degrade at the same rate under parameter perturbation ($\sim$42\% loss at $\sd = 0.5$).

\item \textbf{A local graph-theoretic property predicts viability.} The net flow into the E-coupled species separates alive from dead wirings ($p = 0.0007$). Augmentations survive if and only if the energy-coupling point disperses rather than accumulates concentration from the motif subgraph.
\end{enumerate}

Together with Papers~3--4, these results paint a picture of prebiotic transient complexity as a fragile, topology-sensitive phenomenon. The evolutionary ridge is narrow in parameter space and easily destroyed by network growth---but when the structural constraint at the coupling interface is satisfied, selection can maintain and extend the exploratory phase.

\begin{acknowledgments}
This research used computational resources from Amazon Web Services (EC2 spot instances, total cost $< \$2$). The author thanks Claude (Anthropic), ChatGPT (OpenAI), and Gemini (Google) for extensive collaboration on experimental design, code development, and analysis.
\end{acknowledgments}

\begin{thebibliography}{99}

\bibitem{paper1}
C.~Merrill, ``Dynamical Activation in Autocatalytic Chemical Networks: A Correlation Dimension Analysis,'' Paper~1 in this series (2026).

\bibitem{paper2}
C.~Merrill, ``Feedback-Aligned Autocatalysis Preserves Dynamical Viability but Not Activation in Growing Chemical Networks,'' Paper~2 in this series (2026).

\bibitem{paper3}
C.~Merrill, ``Transient Dimensional Inflation via Timescale Separation in Energy-Coupled Chemical Oscillators,'' Paper~3 in this series (2026).

\bibitem{paper4}
C.~Merrill, ``Evolutionary Selection on Transient Dynamical Retention Time in Energy-Coupled Chemical Oscillators,'' Paper~4 in this series (2026).

\bibitem{feinberg1972}
M.~Feinberg, ``Complex balancing in general kinetic systems,'' Arch.\ Ration.\ Mech.\ Anal.\ \textbf{49}, 187 (1972).

\bibitem{horn1972}
F.~Horn and R.~Jackson, ``General mass action kinetics,'' Arch.\ Ration.\ Mech.\ Anal.\ \textbf{47}, 81 (1972).

\bibitem{nghe2015}
P.~Nghe \textit{et~al.}, ``Prebiotic network evolution: six key parameters,'' Mol.\ BioSyst.\ \textbf{11}, 3206 (2015).

\bibitem{vanchurin2022}
V.~Vanchurin, Y.~I.~Wolf, M.~Katsnelson, and E.~V.~Koonin, ``Toward a theory of evolution as multilevel learning,'' Proc.\ Natl.\ Acad.\ Sci.\ USA \textbf{119}, e2120037119 (2022).

\bibitem{szostak2012}
J.~W.~Szostak, ``The eightfold path to non-enzymatic RNA replication,'' J.\ Syst.\ Chem.\ \textbf{3}, 2 (2012).

\bibitem{researchplan}
C.~Merrill, ``Paper~5 Research Plan: Topological Buffering of Transient Complexity,'' pre-registered experimental design (2026).

\end{thebibliography}

\end{document}
